\subsection{ADJUNCTION OF A UNIT ELEMENT TO A PSEUDO-RING}

Let A be a pseudo-ring (I, § 8, no. 1). On the set \(\mathbf{A}' = \mathbf{Z} \times \mathbf{A}\) we define the following laws of composition:

\[
\left\{ \begin{array}{l} (m, a) + (n, b) = (m + n, a + b) \\ (m, a) (n, b) = (m n, m b + n a + a b). \end{array} \right.\tag{1}
\]

It is immediately verified that A' with these two laws of composition is a ring in which the element \((1,0)\) is the unit element. The set \(\{0\} \times A\) is a two-sided ideal of A' and \(\iota: x \mapsto (0, x)\) is an isomorphism of the pseudo-ring A onto the sub-pseudo-ring \(\{0\} \times A\) by means of which A and \(\{0\} \times A\) are identified. A' is called the ring derived from the pseudo-ring A by adjoining a unit element.

If A already has an identity element \(\varepsilon\) , the element \(e = (0, \varepsilon)\) of \(\mathbf{A}'\) is an idempotent belonging to the centre of \(\mathbf{A}'\) and such that

\[
\mathrm{A} = e \mathrm{A} ^ {\prime} = \mathrm{A} ^ {\prime} e.
\]

Then \((e\mathbf{A}', (1 - e)\mathbf{A}')\) is a direct decomposition (I, § 8, no. 11) of \(\mathbf{A}'\) and the ring \((1 - e)\mathbf{A}'\) is isomorphic to \(\mathbf{Z}\) .

